\documentclass[12 pt]{exam}
\usepackage[margin=0.75in]{geometry}
\usepackage{amsmath,amssymb,color,amsthm}
\usepackage{enumerate,graphicx,multicol}
\usepackage{wrapfig}
\usepackage{pgf,tikz, subcaption, hyperref}
\usetikzlibrary{arrows}
\printanswers %happy last assignment
\newtheorem{claim}{Claim}
\newtheorem{lemma}{Lemma}

\pagestyle{head}                       % This causes the exam to only have headers, no footers.
%\runningheadrule                     % This causes the headings to have an underline.

%\firstpageheader{Name:\enspace\hbox to 4.5in{\hrulefill} Section:\enspace\hbox to 1.2in{\hrulefill}}{}{}
%\extraheadheight{.25in}

\begin{document}
\hrule
\begin{center}
  \huge{Math 2710 Problem Set 14: Ch 14}    %Title
\end{center}
\hrule
\vskip .2in


\begin{questions}

\question Define $f:\mathbb{N} \to \mathbb{Z}$ as $$f(n)=\begin{cases} \frac{n}{2} & n \textrm{ is even } \\ -\frac{n-1}{2} & n \textrm{ is odd }\end{cases}.$$  Show that $f(n)$ is a bijective function.  

\question Show that the cardinality of the even integers is equal to the cardinality of the multiples of 3 by constructing a bijection between the two sets. Explain why your function is a bijection.

\question Prove, using induction, that the union of a finite number of countably infinite sets is countably infinite. 

\question Prove or disprove: There is a bijective function $f:\mathbb{Q} \to \mathbb{R}.$ (Hint: the composition of bijective functions is bijective). 

\question Prove or disprove: The cardinality of the set $[a,b]$ is the same as the cardinality of the set $[0,1]$.



\section*{Challenge Problems} 

\question Review the following proof. Determine the proof method used (or attempted) and give the proof a grade of A (complete, correct, and clear), C (partially correct, partially incomplete, and/or partially unclear), or F (not clear, not correct, or not complete). If you give a grade lower than A, explain your reasoning and describe a way to revise the proof for an A, if possible.

\begin{claim}The sets $(0, \infty)$ and $[0, \infty)$ are numerically equivalent. \end{claim}
\begin{proof}
Define the function $f:(0 , \infty) \to [0, \infty)$ by $f(x)=x.$ First, we show that $f$ is one-to-one. Assume that $f(a)=f(b)$. Then $a=b$ and so $f$ is one-to-one. Next, we show that $f$ is onto. Let $r\in[0,\infty)$. Since $f(r)=r$, the function $f$ is onto. Since $f$ is bijective, $|(0,\infty)|=|[0,\infty)|$. 
\end{proof}

\question Prove the result above using the Schr{\"o}der-Bernstein Theorem, i.e prove the sets $(0, \infty)$ and $[0, \infty)$ are numerically equivalent. 

\question Consider the function $f:[0,1) \to (0,1)$ defined by 
$$f(x)=\begin{cases} 
1/2 & x=0 \\
\frac{1}{2^{n+1}} &  x=\frac{1}{2^n} \quad \textrm{ for } n \in \mathbb{N} \\
x & \textrm{otherwise}
\end{cases}$$
\begin{parts}
\part Explain why $f$ is one-to-one. 
\part Explain why $f$ is onto. 
\part Use this to show that the sets $(0, \infty)$ and $[0, \infty)$ are numerically equivalent. 
\end{parts}


\end{questions}
\end{document}
